% !PRINTED 479
\weberpage{479}
Im Falle 3. setzen wir im System II. des § 73 ($n=3$):
\[
u_1=\sqrt[4]{2},\quad v_1^3=f_1(\sqrt{-18})^3=\sqrt[4]{2}\,x,
\]
\[
A_1=\frac{2}{x^2}-\frac{x^2}{2},\quad B_1=2x+\frac{4}{x},
\]
\[
x^4-4x^3-8x-4=(x^2+2)(x^2-4x-2),
\]
also
\weq{4}{x^2-4x-2=0,\quad f_1(\sqrt{-18})^3=\sqrt[4]{2}\,x.}
Die Auflösung von (4) ergibt:
\weq{5}{x=2+\sqrt6.}

Im Falle 4. ist:
\[
u=f(i)=\sqrt[4]{2},\quad v^3=f(\sqrt{-9})^3=\sqrt[4]{2}\,x,
\]
\[
A=\frac{x^2}{2}+\frac{2}{x^2},\quad B=2x-\frac{4}{x},
\]
\[
x^4-4x^3+8x+4=(x^2-2x-2)^2=0,
\]
also
\weq{6}{x^2-2x-2=0,\quad f(\sqrt{-9})^3=\sqrt[4]{2}\,x,}
\weq{7}{x=1+\sqrt3.}

Im Falle 5.:
\[
u=f(i)=\sqrt[4]{2},\quad v=f(\sqrt{-25})=\sqrt[4]{2}\,x,
\]
\[
A=x^3+\frac{1}{x^3},\quad B=2\Bigl(x^2-\frac{1}{x^2}\Bigr),
\]
\[
x^3+\frac{1}{x^3}-2\Bigl(x^2-\frac{1}{x^2}\Bigr)=\Bigl(x+\frac{1}{x}\Bigr)\Bigl(x-\frac{1}{x}-1\Bigr)^2=0.
\]
\weq{8}{x-\frac{1}{x}-1=0,\quad f(\sqrt{-25})=\sqrt[4]{2}\,x,}
\weq{9}{x=\frac{1+\sqrt5}{2}.}

Endlich setzen wir für $n=7$:
\[
u=\sqrt[4]{2},\quad v=f(\sqrt{-49})=\frac{x}{\sqrt[4]{2}}
\]
und finden
\[
\begin{aligned}
0 &= (x^8-4x^7+28x^4-32x+16)\\
  &= (x^4-2x^3-2x^2-4x+4)^2,
\end{aligned}
\]
woraus leicht durch Auflösung einer quadratischen Gleichung:
\weq{10}{x+\frac{2}{x}=1+\sqrt7,\quad x=\sqrt[4]{2}\,f(\sqrt{-49}).}
